% Part: lambda-calculus
% Chapter: introduction
% Section: computable-lambda

\documentclass[../../../include/open-logic-section]{subfiles}

\begin{document}

\olfileid{lam}{int}{lrp}
\olsection{توابع محاسبه‌پذیر \usetoken{S}{lambda definable} هستند}

\begin{thm}
\ollabel{thm:computable-lambda}
هر تابع جزئی محاسبه‌پذیر !!{lambda definable} است.
\end{thm}

\begin{proof}
باید نشان دهیم که هر تابع جزئی محاسبه‌پذیر~$f$ به‌وسیلهٔ یک ترم
لامبدای~$F$ !!{lambda defined} است. بنا بر قضیهٔ صورت نرمال کلینی، کافی
است نشان دهیم که هر تابع بازگشتی اولیه به‌وسیلهٔ یک ترم لامبدا
!!{lambda defined} است و سپس ثابت کنیم که توابع !!{lambda definable}
تحت ترکیب‌های مناسب و جست‌وجوی نامحدود بسته‌اند. برای نشان‌دادن اینکه هر
تابع بازگشتی اولیه به‌وسیلهٔ یک ترم لامبدا !!{lambda defined} است، کافی
است نشان دهیم که توابع آغازین !!{lambda definable} هستند و توابع جزئی
!!{lambda definable} تحت ترکیب، بازگشت اولیه و جست‌وجوی نامحدود بسته‌اند.
\end{proof}

برای خواناترشدن ادامهٔ اثبات از نمادی متعارف‌تر استفاده می‌کنیم. برای
مثال، $M(x, y, z)$ را به‌جای~$Mxyz$ می‌نویسیم. هرچند این نماد القاکننده
است، باید به یاد داشت که ترم‌های حساب لامبدای بدون نوع، آریتیِ همراه
ندارند؛ ازاین‌رو برای همان ترم~$M$ نوشتن~$M(x,y)$ و~$M(x,y,z,w)$ به یک
اندازه معنادار است. بااین‌حال، این نماد نشان می‌دهد که~$M$ را تابعی از
سه متغیر در نظر گرفته‌ایم و مقصود تعریف‌ها را روشن‌تر می‌کند. به همین
شیوه، می‌گوییم «$M$ را با~$M(x,y,z) = \dots$ تعریف کنید» به‌جای اینکه
بگوییم «$M$ را با~$M =
\lambd[x][\lambd[y][\lambd[z][\dots]]]$ تعریف کنید».

\end{document}
